\section{Introduction}

The minimal-model problem asks whether a smooth projective variety whose
canonical divisor is pseudo-effective admits a birational model with nef
canonical divisor. Here pseudo-effective means that the numerical class lies
in the closure of the cone generated by effective divisors, and nef means
nonnegative degree on every integral curve. The models allowed by the minimal
model program are normal and may have terminal singularities. A minimal
model is required to improve canonical discrepancies and to extract no
divisors, as well as to have nef canonical divisor. These requirements retain
the birational information of the original variety; the precise comparison
used here is part of Theorem~\ref{main:theorem}.

We prove existence in the case where the canonical divisor has numerical
dimension one, using its growth of sections with a fixed ample twist. The
argument begins with minimal model programs for positive boundary
perturbations, but its final step is a criterion on surfaces. In particular,
the proof does not require a general minimal-model existence theorem.

\subsection{The invariant and the main theorem}

For a Cartier divisor \(D\) on a smooth projective \(n\)-fold \(X\),
we use Nakayama's section-growth invariant \cite{Nak04} in the following form:
\begin{equation}\label{main:sigma}
\kappa_\sigma(X,D)=
\max\left\{k\in\{0,\ldots,n\}:
\begin{array}{l}
\text{there is an ample Cartier divisor \(A\) with}\\[2pt]
\displaystyle\limsup_{m\to\infty}
\frac{h^0(X,\mathcal O_X(mD+A))}{m^k}>0
\end{array}\right\}.
\end{equation}
The maximum is \(-\infty\) if the set is empty, and \(m\) ranges over
positive integers. The twist \(A\) is fixed before taking the limsup. Thus
\(\kappa_\sigma(X,D)=1\) rules out positive quadratic limsup for
\emph{every} fixed ample Cartier twist. We will use exactly this consequence,
including when quadratic growth is obtained only along the multiples of one
fixed positive integer. It does not assert an upper bound of the form
\(h^0(X,mD+A)=O(m)\).

This growth convention should be distinguished from the intersection
formula for a nef divisor,
\[
\nu(D)=\max\{k\in\{0,\ldots,n\}:D^k\cdot H^{n-k}>0\},
\]
where \(H\) is ample. We do not use that formula for \(K_X\) before
constructing a nef model, nor do we need a comparison theorem between
numerical-dimension conventions.

\begin{samepage}
\begin{theorem}\label{main:theorem}
Let \(X\) be a smooth connected complex projective variety of dimension
\(n\geq3\). Suppose that \(K_X\) is pseudo-effective and
\(\kappa_\sigma(X,K_X)=1\), with the convention
\eqref{main:sigma}. Then there are a normal projective
\(\mathbb Q\)-factorial terminal variety \(Y\) with \(K_Y\) nef and
a finite sequence
\(\phi:X\dashrightarrow Y\) of \(K\)-negative divisorial contractions
and flips. The inverse of \(\phi\) contracts no prime divisor. On a
common smooth projective resolution \(p:W\to X\), \(q:W\to Y\),
compatible canonical divisors satisfy
\begin{equation}\label{main:comparison}
p^*K_X=q^*K_Y+E,
\end{equation}
where \(E\) is an effective \(q\)-exceptional \(\mathbb Q\)-divisor
whose support contains the strict transform of every prime divisor of
\(X\) contracted by \(\phi\).
\end{theorem}
\end{samepage}

The theorem gives a positive resolution of the numerical-dimension-one
case of the minimal-model conjecture for smooth complex projective varieties
of dimension at least three. In particular it applies to smooth fivefolds,
with no hypothesis on \(\chi(X,\mathcal O_X)\).

A \emph{good} minimal model additionally requires a positive multiple of
\(K_Y\) to be generated by global sections. This is the semiampleness
conclusion of abundance. Nonvanishing asks for a nonzero section of a
positive multiple of the canonical divisor. Neither conclusion is asserted by
Theorem~\ref{main:theorem} alone; the fixed ample twist in
\eqref{main:sigma} is retained throughout the section-growth argument.
The following separate consequence uses log abundance.

\begin{corollary}[Good minimal model]\label{cor:good-model}
Under the hypotheses of Theorem~\ref{main:theorem}, the same
\(\mathbb Q\)-factorial terminal endpoint \(Y\) has semiample \(K_Y\);
hence \(Y\) is a good minimal model of \(X\).
\end{corollary}
\begin{proof}
Apply the log abundance theorem \cite[Theorem~1.1]{OpenAILogAbundance2026}
to the projective lc pair \((Y,0)\) over \(\mathbb C\), whose
\(\mathbb Q\)-Cartier canonical divisor is nef.
\end{proof}

\subsection{The problem and the methods in the literature}

The three-dimensional program developed through work of Reid, Benveniste,
Kawamata, Shokurov, Koll\'ar, Miyaoka and others, with Mori's flip theorem
providing a central step in minimal-model existence \cite{Mori88};
see \cite{KM98} for the broader development.
The minimal model program replaces a variety by successive contractions
and flips along canonical-negative extremal rays; the standard foundations,
including the singularities and discrepancy comparisons used here, are
presented in Koll\'ar--Mori \cite{KM98}. Birkar's formulation for log pairs
\cite[Conjecture~1.1]{Birkar10} places existence of log minimal models
alongside the Mori fiber space alternative. A decisive advance was the work
of Birkar, Cascini, Hacon and McKernan \cite{BCHM10}, which established
minimal-model existence for varieties of general type and finite minimal
model programs with scaling for klt pairs with big boundary. The positive
perturbations in our proof lie within precisely this established range.
Their boundary transforms remain big; they need not remain ample.

At the other end of the numerical spectrum, Nakayama developed the
section-growth invariants and divisorial Zariski decompositions that underlie
numerical approaches to the canonical divisor \cite{Nak04}. Druel proved
termination results for directed programs in numerical dimension zero
\cite[Corollary~3.4]{Druel11}, and Gongyo extended this method to prove
minimal-model existence for \(\mathbb Q\)-factorial dlt pairs of numerical
log Kodaira dimension zero
\cite[Theorem~1.1]{Gongyo11}. The problem here is to pass
from arbitrarily small positive perturbations to the canonical divisor in
the next numerical dimension. A finite program for each positive parameter
does not by itself terminate their concatenation at parameter zero.

Our first reduction also has a specific methodological predecessor.
Alexeev, Hacon and Kawamata use fourth homology to control certain
four-dimensional flips \cite[Lemma~3.1 and its proof]{AHK07}. We use the
same kind of decrease after a discrepancy argument has removed
codimension-two components on the flipped side. In arbitrary dimension
\(n\), the relevant finite-dimensional space is the span of algebraic
\((n-2)\)-cycle classes in degree \(2n-4\). The proof of
Lemma~\ref{mmp:stabilization} gives the needed statement directly; the
four-dimensional result is a precedent for the method, not an imported
termination theorem in all dimensions.

The second reduction combines two established strands of positivity.
Nadel's multiplier-ideal vanishing theorem \cite{Nadel90} supplies the
vanishing used to extend sections from a very general surface while allowing the perturbation to shrink with the section degree;
we use the formulation in Lazarsfeld
\cite[Theorem~9.4.8]{Laz04II}, \cite[Theorem~2.4]{LazMI}.
Surface Riemann--Roch then converts positive square into quadratic section
growth. A different surface, through a point of a hypothetical negative
curve, turns a vanishing-order estimate into a fixed exceptional divisor.
The Hodge index theorem makes its square uniformly negative. The surface
intersection calculations are classical \cite[Chapter~V, Section~1]{Hartshorne77};
our use of them keeps one surface resolution and one exceptional curve fixed
while the rational perturbations vary. This uniformity is what produces a
contradiction.

For comparison with results concerning sections of the unperturbed
canonical divisor, Lazi\'c and Peternell prove nonvanishing for an
\emph{already minimal} projective terminal variety with numerical dimension
one and nonzero Euler characteristic \cite[Theorem~6.7]{LP18}.
Liu and Xu prove existence of a good minimal model for smooth projective
varieties of dimension at most five and numerical dimension at most one,
assuming nonnegative Kodaira dimension \cite[Theorem~1.2]{LiuXu25}.
They also obtain a good minimal model for a smooth projective variety of
dimension at most four with nonzero Euler characteristic and
\(0\leq\kappa_\sigma(X,K_X)\leq1\)
\cite[Corollary~5.2]{LiuXu25}.
These hypotheses and conclusions differ from the existence statement above.
We do not use their nonvanishing or abundance results as proof inputs.

\subsection{The route through the proof}

The first step is to study a sequence of shrinking perturbations on one
fixed model.
Starting with an effective ample rational boundary \(B\), we run finite
minimal model programs for decreasing positive coefficients of its
transform. Every step is \(K\)-negative. Discrepancy and cycle-class
arguments show that only finitely many steps change a locus of codimension
at most two; they do not assert termination of the concatenated programs.
Proposition~\ref{mmp:approximation} therefore supplies a finite canonical
prefix \(X\dashrightarrow Y\) such that multiples of
\[
M_j=K_Y+t_jB_Y,\qquad t_j>0,\quad t_j\longrightarrow0,
\]
are generated outside closed subsets \(Z_j\subset Y\) of codimension
at least three. Here \(B_Y\) is the transform of \(B\); the generated
multiple and the bad set may depend on \(j\).

Two surfaces then show that \(K_Y\) is nef. Fix a very ample divisor
\(H\) on \(Y\). A very general complete intersection surface \(S\)
avoids every \(Z_j\) and the singular locus. The restrictions
\(M_j|_S\) are semiample, so their limit \(K_Y|_S\) is nef.
Lemma~\ref{num:growth} uses multiplier-ideal vanishing to extend sections
from this fixed surface with one fixed twist. Positive square
\(K_Y^2\cdot H^{n-2}\) would consequently give quadratic section growth.
Section~\ref{main:proof-section} transfers that growth to the original
\(X\), where \(\kappa_\sigma(X,K_X)=1\) rules it out. Thus the square
is zero.

This first surface need not meet a curve of negative \(K_Y\)-degree.
To exclude such a curve, including one in the singular locus, take another
complete intersection from \(|H|\) through one of its points, and choose
a surface component through that point. On every component the restriction
of \(M_j\) has nonnegative square; these squares, counted with
multiplicity, sum to \(M_j^2\cdot H^{n-2}\to0\). Thus the square on
the chosen surface also tends to zero.
For all sufficiently small perturbations, an ambient jet estimate on a
resolution of \(Y\) forces sections of Cartier multiples \(kM_j\) to vanish to order at
least \(ck\) at one fixed point above the curve, with \(c>0\) independent
of \(j\) and \(k\). On a fixed resolution of the second surface, this
produces an exceptional fixed part. The Hodge index theorem bounds its
square above by \(-\epsilon k^2\), with \(\epsilon>0\) fixed, while
the residual moving system has nonnegative square. Thus the perturbations
have uniformly positive square on that surface, contradicting its
zero-square limit. This is the nefness
criterion of Proposition~\ref{num:nef}. The finite prefix already supplies
the remaining birational and canonical-comparison assertions of
Theorem~\ref{main:theorem}.

\subsection{Conventions}

All varieties and morphisms are over \(\mathbb C\), and all birational
models in the proof are projective. We use compatible canonical divisors.
Log discrepancies are normalized by
\[
a(F;U,\Delta)=1+\operatorname{coeff}_F
\bigl(K_{\widetilde U}-r^*(K_U+\Delta)\bigr),
\]
where \(r:\widetilde U\to U\) is a resolution carrying the prime
divisor \(F\). Terminality at zero boundary means
\(a(F;U,0)>1\) for every exceptional prime \(F\) over \(U\);
an original prime divisor has log discrepancy \(1\).
For an effective rational boundary, klt means that all log discrepancies
are positive \cite[Chapter~2]{KM98}.

A birational map \emph{extracts no divisors} if its inverse contracts no
prime divisor. Intersection numbers of rational Cartier divisors are defined
after clearing denominators. On an integral surface that need not be normal,
we use Cartier intersections with its fundamental cycle
\cite[Chapter~2]{Fulton98}. Characteristic-zero projective resolutions and
common resolutions are used throughout \cite{Hironaka64}.
